\subsection{EXPONENTIAL MAPPING IN THE LINEAR GROUP}

PROPOSITION 17. Let \(\Delta\) be the set of \(z\in \mathbf{C}\) such that \(-\pi < \mathcal{I}(z) < \pi\) and \(\Delta'\) the set of \(z\in \mathbf{C}\) which are not real \(\leqslant 0\) . Let \(E\) be a complete normable space over \(\mathbf{C}\) and \(A\) (resp. \(A'\) ) the set of \(x\in \mathcal{L}(E)\) whose spectrum \(\mathrm{Sp}(x)\) is contained in \(\Delta\) (resp. \(\Delta'\) ). Then \(A\) (resp. \(A'\) ) is an open subset of \(\mathcal{L}(E)\) (resp. \(\mathbf{GL}(E)\) ) and the mappings \(\exp: A\to A'\) and \(\log: A'\to A\) (Spectral Theories, Chapter I, §4, no. 9) are inverse isomorphisms of analytic manifolds.

This follows from Spectral Theories, Chapter I, § 4, Proposition 10 and no. 9.

THEOREM 6. Let E be a real or complex Hilbert space and U the unitary group of E.

\begin{enumerate}
\def\labelenumi{(\roman{enumi})}
\item
  The set \(\mathbf{H}\) of Hermitian elements of \(\mathcal{L}(\mathbf{E})\) is, with the real normed space structure, a closed vector subspace of \(\mathcal{L}(\mathbf{E})\) admitting a topological supplement.
\item
  The set \(\mathbf{H}'\) of elements \(\geqslant 0\) of \(\mathbf{GL}(\mathbf{E})\) is a real analytic submanifold of \(\mathbf{GL}(\mathbf{E})\) .
\item
  The restriction to H of the mapping exp is an isomorphism of real analytic manifolds of H onto H'.
\item
  The mapping \((h,u)\mapsto (\exp h)u\) of \(\mathbf{H}\times \mathbf{U}\) into \(\mathbf{GL}(\mathbf{E})\) is an isomorphism of real analytic manifolds.
\end{enumerate}

Recall that, if \(x \in \mathcal{L}(\mathrm{E})\) , \(x^*\) denotes the adjoint of \(x\) . Let \(\mathrm{H}_1\) be the set of \(x \in \mathcal{L}(\mathrm{E})\) such that \(x^* = -x\) . The formula \(x = \frac{1}{2}(x + x^*) + \frac{1}{2}(x - x^*)\) proves that, with its real normed space structure, \(\mathcal{L}(\mathrm{E})\) is the topological direct sum of \(\mathrm{H}\) and \(\mathrm{H}_1\) , whence (i).

Suppose that \(\mathbf{K} = \mathbf{C}\) . In the notation of Proposition 17, \(\mathrm{H}'\) is the set of \(h \in \mathrm{H} \cap \Delta'\) such that \(\mathrm{Sp} h \subset \mathbf{R}_{+}^{*}\) . As \(\exp(\mathbf{R}) = \mathbf{R}_{+}^{*}\) , (ii) and (iii) follow from Proposition 17 and Spectral Theories, Chapter I, §4, Proposition 8 and §6, no. 5. The mapping \((h, u) \mapsto y = (\exp h)u\) of \(\mathrm{H} \times \mathrm{U}\) into \(\mathbf{GL}(\mathrm{E})\) is bijective by Spectral Theories, Chapter I, §6, Proposition 15. It is real analytic by the above. The mapping \(y \mapsto h = \frac{1}{2} \log(yy^*)\) is real analytic and so also is the mapping \(y \mapsto u = (\exp h)^{-1}y\) . Hence (iv).

Suppose that \(\mathbf{K} = \mathbf{R}\) . Let \(\hat{\mathbf{E}}\) be the complexified Hilbert space of E and J the mapping \(\xi + i\eta \mapsto \xi - i\eta\) (for \(\xi, \eta\) in E) of \(\hat{\mathbf{E}}\) into \(\hat{\mathbf{E}}\) . Let \(\tilde{\mathbf{H}}, \tilde{\mathbf{H}}', \tilde{\mathbf{U}}\) denote the sets defined for \(\hat{\mathbf{E}}\) as were H, H', U for E. Then H (resp. H', U) is identified with the set of \(x \in \tilde{\mathbf{H}}\) (resp. \(\tilde{\mathbf{H}}'\) , \(\tilde{\mathbf{U}}\) ) such that \(\mathrm{JxJ}^{-1} = x\) . Properties (ii), (iii) and (iv) then follow easily from (i) and the analogous properties in the complex case.

PROPOSITION 18. Let E be a complete normable space over C, \(v \in \mathcal{L}(E)\) and g = exp v. Suppose that \(\operatorname{Sp}(v)\) contains none of the points \(2i\pi n\) with \(n \in Z - \{0\}\) . Then, for all \(x \in E\) , the conditions vx = 0 and gx = x are equivalent.

This follows from Lemma 2 of no. 4, applied to the function \(z \mapsto e^z - 1\) .

COROLLARY 1. Let E be a complete normable space over C and F the space of continuous n-linear mappings of \(\mathbf{E}^n\) into E. For all \(v \in \mathcal{L}(\mathbf{E})\) , let \(\sigma(v)\) be the element of \(\mathcal{L}(\mathbf{F})\) defined by

\[
(\sigma (v) f) (x _ {1}, \dots , x _ {n}) = v (f (x _ {1}, \dots , x _ {n})) - \sum_ {i = 1} ^ {n} f (x _ {1}, \dots , v x _ {i}, \dots , x _ {n}).
\]

For all \(g \in \mathbf{GL}(\mathbf{E})\) , let \(\rho(g)\) be the element of \(\mathbf{GL}(\mathbf{F})\) defined by

\[
(\rho (g) f) (x _ {1}, \dots , x _ {n}) = g (f (g ^ {- 1} x _ {1}, \dots , g ^ {- 1} x _ {n})).
\]

Let \(u \in \mathcal{L}(\mathrm{E})\) be such that every \(z \in \operatorname{Sp} u\) satisfies \(|\mathcal{I}(z)| < \frac{2\pi}{n + 1}\) . Then, for all \(f \in \mathrm{F}\) , the conditions \(\sigma(u)f = 0\) and \(\rho(\exp u)f = f\) are equivalent.

\(L(\rho) = \sigma (\S 3, \text{no. } 11, \text{ Corollary 1 to Proposition 41}) \text{ and hence}\)

\[
\rho (\exp u) = \exp \sigma (u)
\]

(no. 4, Corollary 3 to Proposition 10). By Proposition 18 it then suffices to prove that \(\operatorname{Sp} \sigma(u)\) does not meet \(2i\pi(\mathbf{Z} - \{0\})\) . But this follows from the following lemma:

Lemma 6. If \(v \in \mathcal{L}(\mathrm{E})\) , then \(\operatorname{Sp} \sigma(v) \subset \operatorname{Sp} v + \operatorname{Sp} v + \cdots + \operatorname{Sp} v\) , where the sum comprises \(n + 1\) terms.

We define elements \(v_{0}, v_{1}, \ldots, v_{n}\) of \(\mathcal{L}(\mathrm{F})\) by writing, for all \(f \in \mathrm{F}\) ,

\[
\begin{array}{l} (v _ {0} f) (x _ {1}, \dots , x _ {n}) = v (f (x _ {1}, \dots , x _ {n})) \\ (v _ {i} f) (x _ {1}, \dots , x _ {n}) = - f (x _ {1}, \dots , v x _ {i}, \dots , x _ {n}) \quad \text { for } 1 \leqslant i \leqslant n. \end{array}
\]

Then \(\sigma(v) = \sum_{i=0}^{n} v_i\) and the \(v_i\) are pairwise permutable. Let A be the total closed subalgebra of \(\mathcal{L}(F)\) generated by the \(v_i\) ; it is commutative (Spectral

Theories, Chapter I, § 1, no. 4) and \(\mathrm{Sp}_{\mathcal{L}(\mathbb{F})}v' = \mathrm{Sp}_A v' \subset \sum_{i=0}^{n} \mathrm{Sp} v_i\) (Spectral Theories, Chapter I, § 3, Proposition 3 (ii)). Now, if \(\lambda \in \mathbf{C}\) is such that \(v - \lambda\) is invertible, clearly the \(v_i - \lambda_i\) are invertible and hence \(\mathrm{Sp} v_i \subset \mathrm{Sp} v\) for all \(i\) .

COROLLARY 2. Let \(\mathbf{E}\) be a complete normable algebra over \(\mathbf{C}\) and \(w \in \mathcal{L}(\mathbf{E})\) . Suppose that every \(z \in \operatorname{Sp} w\) satisfies \(|\mathcal{I}(z)| < \frac{2\pi}{3}\) . The following conditions are equivalent:

\begin{enumerate}
\def\labelenumi{(\roman{enumi})}
\item
  \(w\) is a derivation of \(\mathbf{E}\) ;
\item
  exp w is an automorphism of E.
\end{enumerate}

This follows from Corollary 1 with \(n = 2\) and \(f\) the multiplication of E.

PROPOSITION 19. Let E be a complete normable space over C, \(v \in \mathcal{L}(E)\) and g = exp v. Suppose that every \(z \in Sp\) v satisfies \(-\pi < \mathcal{I}(z) < \pi\) . Then, for every closed vector subspace \(E'\) of E, the conditions \(v(E') \subset E'\) and \(g(E') = E'\) are equivalent.

The condition \(v(\mathrm{E}') \subset \mathrm{E}'\) implies \(g(\mathrm{E}') \subset \mathrm{E}'\) and \(g^{-1}(\mathrm{E}') \subset \mathrm{E}'\) and hence \(g(\mathrm{E}') = \mathrm{E}'\) . Suppose that \(g(\mathrm{E}') = \mathrm{E}'\) . We use the notation \(\Delta, \Delta'\) of Proposition 17. Since Sp \(v\) is a compact subset of \(\Delta\) , there exists a compact rectangle \(\mathbf{Q} = [\underline{a}, b] \times [\underline{a}', b']\) such that \(\operatorname{Sp} v \subset \mathbf{Q} \subset \Delta\) . The set \(\Delta - \mathbf{Q}\) is connected. Hence \(\operatorname{Sp} g \subset \exp \mathbf{Q} \subset \Delta'\) , the set \(\exp \mathbf{Q}\) is compact and the set \(\Delta' - \exp \mathbf{Q}\) is connected. The closure of the latter contains \(] - \infty, [0]\) and hence

\[
\left. \left(\Delta^ {\prime} - \exp Q\right) \cup ] - \infty , 0 ] = C - \exp Q \right.
\]

is connected. Then exp Q is polynomially convex (Spectral Theories, Chapter I, §3, Corollary 2 to Proposition 9) and hence the function log, defined on \(\Delta'\) , is a limit in \(\mathcal{O}(\exp Q)\) of polynomial functions (Spectral Theories, Chapter I, §4, Proposition 3). Hence \(v = \log g\) is the limit in \(\mathcal{L}(E)\) of elements of the form P(g), where P is a polynomial (Spectral Theories, Chapter I, §4, Theorem 3). As P(g)(E') \(\subset\) E', it follows that \(v(E') \subset E'\) .

COROLLARY. Let E be a complete normed space over C, \(v \in \mathcal{L}(E)\) and \(g = \exp v\) .

Suppose that every \(z \in \operatorname{Sp} v\) satisfies \(-\frac{\pi}{2} < \mathcal{I}(z) < \frac{\pi}{2}\) . Then, for every closed vector

subspace M of \(\mathcal{L}(\mathrm{E})\) , the conditions \(g\mathrm{M}g^{-1} = \mathrm{M}\) and \([v, \mathrm{M}] \subset \mathrm{M}\) are equivalent.

Let \(\mathrm{F} = \mathcal{L}(\mathrm{E})\) , \(g'\) be the mapping \(f \mapsto gfg^{-1}\) of \(\mathrm{F}\) into \(\mathrm{F}\) and \(v'\) be the mapping \(f \mapsto [v, f]\) of \(\mathrm{F}\) into \(\mathrm{F}\) . Then \(g' = \exp v'\) (no. 4, Corollary 3 to Proposition 10 and §3, no. 11, Corollary 1 to Proposition 41). Lemma 6 proves that \(-\pi < \mathcal{I}(z) < \pi\) for all \(z \in Sp v'\) . It then suffices to apply Proposition 19.

